\subsection{Anillos estrictamente formalmente catenarios.}
\label{subsection:IV.7.2-fr}

\begin{env}[7.2.1]
\label{IV.2.s7.2.1-fr}
\emph{Notaciones.} --- Para un anillo \emph{íntegro noetheriano} $A$, denotaremos
por $A^{(1)}$ la intersección de los anillos locales $A_{\mathfrak p}$,
cuando $\mathfrak p$ recorre el conjunto de los ideales primos de $A$ \emph{de
altura} 1 (5.10.17). Si $A$ es un anillo \emph{local} noetheriano íntegro,
denotaremos por $A^{(\omega)}$ la intersección de los $A_{\mathfrak p}$,
cuando $\mathfrak p$ recorre esta vez el conjunto de todos los ideales
primos de $A$ \emph{distintos del ideal maximal}.

Diremos que un anillo local noetheriano $A$ es \emph{estrictamente
equidimensional} si, para todo ideal primo
$\mathfrak p\in\operatorname{Ass}(A)$, se tiene
$\dim(A/\mathfrak p)=\dim(A)$; equivale a decir que $A$ es
\emph{equidimensional} y \emph{carece de ideales primos asociados inmersos}.
\end{env}
